\section{Summary of open questions}
\label{sec:summary-of-open-questions}

\subsection{Open problems in mechanistic interpretability methods and foundations}

\subsubsection{Reverse engineering: Identifying the roles of network components}

\parasection{Reverse engineering step 1: Neural network decomposition}

\begin{enumerate}
    \item How should we decompose networks into more interpretable constituent parts?
    \begin{enumerate}[label=\alph*.]
        \item What isomorphism or what approximation of a neural network (or parts of it) is the best way to express it for the purposes of interpreting it?
        \item How should we coarse grain neural networks?
        \item How should we build higher level abstractions on top of low-level network components?
    \end{enumerate}
    
    \item How true is the linear representation hypothesis?
    \begin{enumerate}[label=\alph*.]
        \item To what extent do models encode concepts linearly in their representations?
        \item How should we characterize representations that are not linearly represented in neural networks?
        \item What properties of a concept, or of the training distribution, result in a particular concept becoming encoded linearly (or not)?
    \end{enumerate}
    
    \item Is the combination of the linear representation hypothesis and superposition the right frame for thinking about computation in neural networks?
    \begin{enumerate}[label=\alph*.]
        \item Can we fully determine the causes of feature superposition and polysemanticity within neural networks?
        \item How should we understand superposition in attention blocks?
        \item How should we understand cross-layer superposition?
        \item What new theoretical insights can be gleaned from considering how networks perform computation natively in superposition, rather than treating superposition purely as a compression strategy?
    \end{enumerate}
    
    \item Can the problems with SDL be overcome?
    \begin{enumerate}[label=\alph*.]
        \item What lies in SDL reconstruction errors? Will the errors converge to zero with methodological progress?
        \item Is sparsity the correct proxy for interpretability?
        \item Can the approach be scaled to the largest models?
        \item Does SDL make sense if we don't believe in the linear representation hypothesis?
        \item Is sparsity the best possible proxy for interpretability?
        \item Is it correct to think of SDL features as `bags of features', or is there important information contained within the geometry of representation space?
        \item If SDL finds compositions of the “true” features, is this a problem?
        \item Can SDL be applied to all architectural components to fully decompose networks?
        \item How can we better measure the success of SDL techniques?
        \item How should we connect sparsely activating features into circuits? Will this be practically possible, and the best possible description of network mechanisms?
        \item Can we develop new methods that address conceptual and practical issues with SDL?
    \end{enumerate}
    
    \item How important is the geometry of activation space for explaining neural network behavior?
    \begin{enumerate}[label=\alph*.]
        \item How can we identify the underlying functional structure of networks (which defines why activations are located in particular geometric arrangements in activation space)?
        \item Must we understand global feature geometry or only local feature geometry in order to understand computation in neural networks?
    \end{enumerate}
    
    \item Can we connect theories for how neural networks generalize to interpretability?
    \begin{enumerate}[label=\alph*.]
        \item Can we distinguish parts of networks that underlie generalization from parts that underlie memorization?
        \item What mechanisms underlie the relationship between interpretability and generalization?
        \item Are there connections between adversarial robustness and superposition?
        \item Can we connect interpretability to theories of deep learning like SLT?
    \end{enumerate}
    
    \item Can we build intrinsically more interpretable models at low performance cost? How helpful is this?
    \begin{enumerate}[label=\alph*.]
        \item Interpretability training: Can we train networks that are interpretable by default at low performance cost?
        \item Interpretable inference: Can we convert already-trained models into forms that are much easier to completely interpret at little performance cost?
        \item How can we train large-scale models such that the concepts they use are naturally understandable to humans?
        \item How can we design training objectives so that the model is incentivized to use known specific abstractions?
        \item How can we localize concepts we want to control (e.g., long-term plans) during training?
    \end{enumerate}
\end{enumerate}

\parasection{Reverse engineering step 2: Identifying the functional role of components}
\begin{enumerate}
    \item Can we improve on max-activating input data set examples for understanding the causes of network component activations?
    \begin{enumerate}[label=\alph*.]
        \item How can we avoid imposing human bias to explanations?
        \item Can we progress toward deeper descriptions based on internal mechanisms?
        \item How might we develop interpretation methods that can recognize and work with unfamiliar concepts - computational patterns that don't map cleanly to human intuitions?
    \end{enumerate}
    
    \item How can we develop attribution methods that faithfully and efficiently compute which network components are important for some downstream metric?
    \begin{enumerate}[label=\alph*.]
        \item How can we develop attribution methods that capture higher-order effects beyond first-order approximations of model behavior?
        \item Is it possible to create perturbation-based methods that don't force models to operate outside their training distribution?
        \item Can we develop hybrid approaches that combine the strengths of different attribution methods while mitigating their individual weaknesses?
    \end{enumerate}
    
    \item How can we better measure the downstream effects of model components?
    \begin{enumerate}[label=\alph*.]
        \item How can we reliably distinguish between true causal pathways and compensatory effects like the “Hydra effect” when performing interventions?
    \end{enumerate}
\end{enumerate}


\parasection{Reverse engineering step 3: Validation of descriptions}
\begin{enumerate}
    \item Can we improve our ability to validate mechanistic explanations for model behavior in ways that do not depend on researcher intuition and are computationally tractable to use?
    \begin{enumerate}[label=\alph*.]
        \item Can we improve on methodologies for evaluating hypotheses through their predictive power on activations of network components?
        \item Can we develop methodologies for evaluating hypotheses through their predictive power for model behavior (e.g. unusual failures or adv examples)?
        \item Can we develop methodologies for handcoding weights that faithfully represent some hypotheses, as a drop in replacement for subnetworks we claim to understand?
        \item Can we develop a wider suite of networks with known ground truth explanations to validate techniques against?
        \item Can we use mechanistic explanations to achieve engineering goals?
        \item Can we use mechanistic explanations to achieve engineering goals in a way that improves upon black box baselines?
    \end{enumerate}
    
    \item Can we develop “model organisms” as a community, which are understood deeply, and seen as a test-bed for new unproven interpretability methodologies to be tested?
    
    \item Can we establish standardized baselines and benchmarks for comparing different interpretability approaches on real-world, non-cherry-picked tasks, where the ground truth is known?
    
    \item What would constitute a comprehensive set of “stress tests” for interpretability hypotheses that could reliably detect interpretability illusions?
    
    \item How might we design evaluation frameworks that assess interpretability methods on their average case and worst-case performance rather than just best-case scenarios?
    
    \item How can we ensure that our understanding of internals generalizes to out-of-distribution inputs?
\end{enumerate}

\subsubsection{Concept-based interpretability: Identifying network components for given roles}
\begin{enumerate}
    \item How can we reliably distinguish causal from merely correlated features when probing neural networks?
    
    \item Can we develop automated systems to generate high-quality probing data sets, reducing the current heavy reliance on human effort?
    
    \item What regularization and validation techniques can be used to prevent spurious correlations while ensuring probes find generalizable features?
    
    \item How can we improve probing for concepts that may not have clear positive/negative examples?
\end{enumerate}

\subsubsection{Proceduralizing mechanistic interpretability into circuit discovery pipelines}
\begin{enumerate}
    \item Can we develop techniques that build on lower level methods that provide deeper or more complete insights about neural networks?
    \item How much can we learn from further work in the existing circuit discovery paradigm?
    \begin{enumerate}[label=\alph*.]
        \item Should we expect circuit discovery to benefit from further methodological progress in decomposing neural networks? Will faithfulness go up and explanation description length go down?
        \item Can we remove the constraint that task definition for circuit discovery is inherently concept-based, which may be complicating mechanistic analysis?
        \item Can we get around the practical issues of negative and backup behavior?
        \item Will circuit discovery provide insights into arbitrary tasks, or will it only be helpful in cases where we are able to crisply define tasks?
    \end{enumerate}
\end{enumerate}

\subsubsection{Automating steps in mechanistic interpretability research}
\begin{enumerate}
    \item Can we improve on AI automated feature description and validation methods?
    \begin{enumerate}[label=\alph*.]
        \item Through automating the generation and testing of arbitrary hypotheses?
        \item Through describing differences between features?
        \item Through descriptions of how components interact?
    \end{enumerate}
    
    \item Can we improve on ACDC-like circuit discovery methods?
    
    \item Can we automate other parts of the mechanistic interpretability pipeline?
    \begin{enumerate}[label=\alph*.]
        \item Conceptual interpretability research?
        \item Decomposition method discovery?
        \item More ad hoc validation of hypotheses?
    \end{enumerate}
    
    \item Should we take steps to mitigate potentially misaligned AI systems sabotaging AI automated interpretability?
\end{enumerate}

\subsection{Open problems in applications of mechanistic interpretability}

\subsubsection{Using mechanistic interpretability for better monitoring and auditing of AI systems for potentially unsafe cognition}
\begin{enumerate}
    \item Can we effectively use interpretability for safety evaluations?
    \begin{enumerate}[label=\alph*.]
        \item Can we develop robust “white box” evaluations that detect concerning internal patterns without needing to understand the entire network?
        \item Can we reliably distinguish between features that merely recognize deceptive behavior versus mechanisms that generate deceptive behavior?
        \item How can we validate that learned features capture all concerning patterns of reasoning?
        \item Can we reliably identify which features are appropriately versus spuriously relevant to a given task?
    \end{enumerate}
    
    \item Can we leverage interpretability to enhance red-teaming and system testing?
    \begin{enumerate}[label=\alph*.]
        \item Can we use interpretability insights to make red-teaming more efficient than current methods?
        \item How can we best use feature attribution to help human red-teamers identify problematic input patterns?
    \end{enumerate}
    
    \item Can we develop effective test-time monitoring systems based on interpretability?
    \begin{enumerate}[label=\alph*.]
        \item Can we get mechanistic anomaly detection to work?
        \item Can we create passive monitoring systems based on model internals that effectively flag concerning internal patterns during deployment?
        \item Can we develop monitoring systems that work with only feature-level understanding rather than requiring deep mechanical insights?
    \end{enumerate}
\end{enumerate}

\subsubsection{Using mechanistic interpretability for better control of AI system behavior}
\begin{enumerate}
    \item Can we improve steering methods through interpretability?
    \begin{enumerate}[label=\alph*.]
        \item How can we make activation steering more precise and reduce its side effects?
        \item Can we develop methods to steer entire mechanisms rather than just single features?
    \end{enumerate}
    
    \item Can we achieve reliable model unlearning and editing?
    \begin{enumerate}[label=\alph*.]
        \item Will carving the network at its true joints help us improve on model unlearning and editing?
        \item Can mechanistic interpretability help us develop better methods for evaluating unlearning efficacy?
        \item Can mechanistic interpretability help us determine which classes of model edit are even possible, without damaging generalization in undesirable ways?
    \end{enumerate}
    
    \item Can we better understand and improve finetuning through interpretability?
    \begin{enumerate}[label=\alph*.]
        \item Can we make finetuning more sample-efficient by targeting specific parameters?
        \item Can we develop better tools for analyzing feature-level or mechanism-level differences between model versions?
    \end{enumerate}
\end{enumerate}

\subsubsection{Using mechanistic interpretability for better predictions about AI systems}
\begin{enumerate}
    \item Can we predict model behavior in novel situations outside of the distribution of inputs we have access to with mechanistic understanding?
    \begin{enumerate}[label=\alph*.]
        \item Can we reliably predict when and how jailbreaking or safety bypasses might occur?
        \item How can we identify internal signatures that predict specific failure modes like hallucination?
        \item Can we develop methods to predict model behavior without requiring behavioral evaluations?
        \item Can we find “values” or “goals” in systems that might be indicative of behavior in more generality?
        \item Is it possible to prove the absence of specific dangerous capabilities through mechanistic analysis?
    \end{enumerate}
    
    \item Can we develop formal verification methods for AI systems?
    \begin{enumerate}[label=\alph*.]
        \item Can current toy model verification approaches scale to frontier systems?
        \item How much of neural computation can be reduced to verifiable symbolic operations?
        \item Can we create formal guarantees about system behavior in complex, non-formalizable environments?
        \item What level of mechanistic understanding is necessary for meaningful formal verification?
    \end{enumerate}
    
    \item Can we make rigorous claims about model safety?
    \begin{enumerate}[label=\alph*.]
        \item Can we definitively prove the absence of specific dangerous mechanisms?
        \item How can we verify claims about model values and goals in a rigorous way?
        \item What types of safety claims are possible with current interpretability methods?
        \item Can we develop “enumerative safety” approaches that reliably identify all relevant mechanisms?
    \end{enumerate}
    
    \item Can we better predict AI capability development through interpretability?
    \begin{enumerate}[label=\alph*.]
        \item Can we identify early signatures that predict emergent capabilities?
        \item How do model mechanisms evolve dynamically during training?
        \item Can we map the connection between small-scale circuits and large-scale capabilities?
        \item How does the loss landscape's structure relate to capability emergence?
    \end{enumerate}
    
    \item Can we understand the relationship between training data and capabilities?
    \begin{enumerate}[label=\alph*.]
        \item How do specific training examples influence the development of model mechanisms?
        \item Can we predict model limitations based on training data composition?
        \item Can we design training data sets to reliably produce specific desired capabilities?
        \item How does data set structure affect the balance between in-context and weights-based learning?
    \end{enumerate}
    
    \item Can we predict latent or maskable capabilities?
    \begin{enumerate}[label=\alph*.]
        \item How can we identify capabilities that could be `unlocked' through prompting or finetuning?
        \item Can we detect when finetuning has masked rather than removed capabilities?
        \item How do we analyze mechanisms that span multiple timesteps or sequential behaviors?
        \item Can we predict which model capabilities are fundamental versus superficially trained?
    \end{enumerate}
\end{enumerate}

\subsubsection{Using mechanistic interpretability to improve our ability to perform inference, improve training and make use of learned representations}
\begin{enumerate}
    \item Can we use interpretability to make inference more efficient?
    \begin{enumerate}[label=\alph*.]
        \item How can we identify skippable computations without affecting outputs?
        \item Can we create more effective distillation methods through mechanistic understanding?
        \item How can we optimize model architecture based on component function analysis?
        \item Can we identify and optimize critical computational pathways?
    \end{enumerate}
    
    \item Can we improve training through mechanistic insights?
    \begin{enumerate}[label=\alph*.]
        \item Can we better select training data by understanding example influence?
        \item How can we monitor and optimize capability emergence during training?
        \item Can we develop more parameter-efficient training methods through component analysis?
        \item Can we create better architectures through component understanding?
        \item Can we identify and enhance components with specific functionalities?
    \end{enumerate}
    
    \item Can we instill capabilities directly into networks?
    \begin{enumerate}[label=\alph*.]
        \item Can we design better inductive biases based on mechanistic insights?
        \item Is it possible to create modular architectures with swappable components?
        \item Can we develop reliable methods for combining model parameters?
        \item Is it possible to transfer specific capabilities between models?
    \end{enumerate}
\end{enumerate}

\subsubsection{Using mechanistic interpretability for 'microscope AI'}
\begin{enumerate}
    \item Can we leverage AI models for scientific discovery?
    \begin{enumerate}[label=\alph*.]
        \item How can we extract novel patterns and predictors that models have found?
        \item Can we make microscope AI techniques accessible to domain experts?
        \item How do we validate scientific insights derived from model interpretability?
        \item Can we extend microscope AI beyond current simple correlational discoveries?
    \end{enumerate}
    
    \item Can we develop better knowledge extraction methods?
    \begin{enumerate}[label=\alph*.]
        \item How can we detect when models have found genuinely novel patterns?
        \item Can we automate the process of finding scientific insights in model weights?
        \item How do we bridge the gap between model features and scientific concepts?
        \item Can we make these techniques usable without deep machine learning expertise?
    \end{enumerate}
\end{enumerate}

\subsubsection{Mechanistic interpretability on a broader range of models and model families}
\begin{enumerate}
    \item Can interpretability methods generalize across architectures?
    \begin{enumerate}[label=\alph*.]
        \item Do current interpretability methods (SDL, circuit analysis) transfer to SSMs? Or, like the transition from CNNs to transformers, are new approaches necessary?
        \item Which insights are model-specific versus universal?
        \item How can we adapt methods for multimodal models?
    \end{enumerate}
    
    \item How do different models trained on similar data compare mechanistically?
    \begin{enumerate}[label=\alph*.]
        \item Is the “universality hypothesis” true across models? To what extent do neural networks learn similar features and circuits to each other (and to humans?)
        \item Do different architectures learn fundamentally different features?
        \item How do mechanisms of particular tasks differ between transformers, CNNs, and SSMs?
        \item Are there insights we can gain from comparing architectures?
    \end{enumerate}
    
    \item Can we future-proof interpretability research?
    \begin{enumerate}[label=\alph*.]
        \item How can we prepare for interpreting novel architectures?
        \item Should we focus on architecture-specific or general methods?
        \item Can we identify truly fundamental interpretability principles?
        \item Will current methods work on future frontier models?
    \end{enumerate}
\end{enumerate}

\subsubsection{Human computer interaction with model internals}
\begin{enumerate}
    \item Can we create interfaces that use mechanistic understanding to enhance human-neural network interaction?
    \begin{enumerate}[label=\alph*.]
        \item How can we visualize model internals in an intuitive way?
        \item Can we develop real-time interpretability dashboards?
        \item What's the right balance between simplicity and depth in these interfaces?
        \item How do we make complex model mechanisms understandable to non-experts?
    \end{enumerate}
    
    \item Can we develop interpretability tools to help auditors?
    \begin{enumerate}[label=\alph*.]
        \item How can we help auditors find potential failure modes directly?
        \item Can we develop tools to detect bias at the mechanism level?
        \item What interfaces would make auditing more efficient and thorough?
        \item How can we present technical findings to policy makers?
    \end{enumerate}
    
    \item Can we improve end-user interaction with AI?
    \begin{enumerate}[label=\alph*.]
        \item How can transparency features help users calibrate trust?
        \item Can we create intuitive controls based on model mechanisms?
        \item Can we create intuitive ways to steer model behavior?
    \end{enumerate}
\end{enumerate}

\subsubsection{Governance}
\begin{enumerate}
    \item Can mechanistic analysis help identify and prevent failures?
    \begin{enumerate}[label=\alph*.]
        \item Can we identify specific mechanisms that caused AI failures?
        \item How do we map the causal chain of mechanisms leading to incidents?
        \item Can we detect when similar mechanisms are about to activate?
        \item Is it possible to isolate and modify failure-causing mechanisms?
    \end{enumerate}
    
    \item Can we study mechanism patterns related to governance?
    \begin{enumerate}[label=\alph*.]
        \item Can we identify mechanisms responsible for specific dangerous capabilities?
        \item How do we detect deceptive or evasive mechanisms?
        \item Can we map the mechanisms involved in model decision-making?
        \item Is it possible to verify the absence of specific harmful mechanisms?
    \end{enumerate}
    
    \item Can mechanistic insights verify compliance?
    \begin{enumerate}[label=\alph*.]
        \item How can we trace decision mechanisms to explain model outputs?
        \item Can we identify mechanisms that process copyrighted content?
        \item Is it possible to detect mechanisms that encode specific knowledge?
        \item How do we verify modifications to problematic mechanisms?
    \end{enumerate}
\end{enumerate}

\subsubsection{Open socio-technical problems in mechanistic interpretability}



\parasection{Translating technical progress in mechanistic interpretability into levers for AI policy and governance}
\begin{enumerate}
    \item Can we use a mechanistic understanding to better evaluate AI capabilities?
    \begin{enumerate}[label=\alph*.]
        \item How can we use interpretability to improve capability elicitation?
        \item Can we use interpretability to reliably detect when models are strategically underperforming capabilities evaluations?
    \end{enumerate}
    
    \item Can we use a mechanistic understanding to improve our ability to forecast when or whether new capabilities will arise ahead of time?

    
    \item How can we use interpretability to better estimate the likelihood of different threat models?


    \item Can we use interpretability to prevent AI incidents?
    \begin{enumerate}[label=\alph*.]
        \item Can we use interpretability to construct reliable test-time monitors to detect AI incidents?
        \item Can we use reliably prevent similar incidents in the future, by using interpretability to design new evaluation tasks on incident scenarios?
    \end{enumerate}
    \item Can interpretability help verify which workloads GPUs are being used for?
    \item How should interpretability inform copyright law?
    \item How can mechanistic understanding help resolve copyright challenges in generative AI, particularly regarding the detection and removal of memorized copyrighted works?
\end{enumerate}


\parasection{Social and philosophical challenges in mechanistic interpretability}
\begin{enumerate}
    \item What is interpretability?
    \begin{enumerate}[label=\alph*.]
        \item What are the goals of the field?
        \item How should success be graded?
        \item Should we treat interpretability as a science or an engineering discipline? What implications does this have on what research should be done?
    \end{enumerate}
    
    \item How can we mitigate downside risks of interpretability research?
    \begin{enumerate}[label=\alph*.]
        \item How can we communicate the results of our research such that the risk of their misuse is minimized?
    \end{enumerate}
\end{enumerate}